\section{Introduction}\label{sec:introduction}

Repulsive interactions move particles out of a Bose condensate even at
zero temperature. For a dilute three-dimensional gas, Bogoliubov theory
predicts a universal leading depletion, depending on the interaction
only through its scattering length. We establish this prediction for
fixed bounded radial repulsive potentials of finite range. The volume
tends to infinity at each fixed density before the density tends to zero.

\subsection{The model and the result}

We use units in which $\hbar^2/(2m)=1$. Let
$v:\R^3\to[0,\infty)$ be bounded, measurable, radial, and compactly
supported, and suppose that $v$ is nonzero on a set of positive measure.
Its scattering length $a_v>0$ is determined by the radial solution of
\[
 (-\Delta+\tfrac12v)f=0,
 \qquad f(x)=1-\frac{a_v}{|x|}
 \quad\text{for sufficiently large }|x|.
\]
On the torus $\Lambda_L=(\R/L\Z)^3$, set
\[
 v_L(x)=\sum_{z\in\Z^3}v(x+Lz),
 \qquad
 H^v_{N,L}=\sum_{j=1}^N(-\Delta_j)
                   +\sum_{i<j}v_L(x_i-x_j),
\]
acting on $L^2_{\mathrm{sym}}(\Lambda_L^N)$ with periodic boundary
conditions. Write $\mathcal G^v_{N,L}$ for the positive trace-one
operators supported on its lowest eigenspace. For
$\Gamma\in\mathcal G^v_{N,L}$, define
\[
 \gamma^{(1)}_\Gamma=N\Tr_{2,\ldots,N}\Gamma,
 \qquad u_0=L^{-3/2},
 \qquad
 B_\Gamma=\frac{\langle u_0,\gamma^{(1)}_\Gamma u_0\rangle}{N}.
\]
Thus $B_\Gamma$ is the occupation fraction of the constant orbital and
$1-B_\Gamma$ is its depletion.

\begin{theorem}\label{thm:main}
For every potential $v$ satisfying the assumptions above and every
$\epsilon>0$, there is $\rho_0(v,\epsilon)>0$ such that, for every fixed
$0<\rho<\rho_0(v,\epsilon)$ and every sequence
$N_k,L_k\to\infty$ with $N_k/L_k^3\to\rho$,
\[
 \limsup_{k\to\infty}\ \sup_{\Gamma\in\mathcal G^v_{N_k,L_k}}
 \left|\frac{1-B_\Gamma}{\sqrt{\rho a_v^3}}
                       -\frac8{3\sqrt\pi}\right|\le\epsilon.
\]
\end{theorem}

The potential remains fixed throughout both limits. In particular, the
statement applies to every thermodynamic accumulation value of the
occupation; existence of an occupation limit at each positive density
is not required. Neither monotonicity of the radial potential nor
positivity of its Fourier transform is assumed. Boundedness makes the
finite-volume heat semigroup positivity improving, so the ground state
is in fact unique. The formulation with density matrices records the
occupation observable without a choice of ground-state vector.

\subsection{Historical context and antecedents}

Bogoliubov's quasiparticle theory describes the excitations and depletion
of a weakly interacting condensate~\cite{Bogoliubov}. Lee, Huang, and
Yang calculated the leading depletion and the second-order energy of
a dilute hard-sphere gas~\cite{LHY}. The two predictions concern
different features of the many-particle ground state: its energy and
its one-body density matrix. In the ordinary thermodynamic limit,
long-wavelength particles can have arbitrarily small kinetic energy.
Consequently, an accurate energy expansion does not by itself determine
the occupation of the constant orbital.

The rigorous energy theory began with Dyson's upper and lower bounds
for hard spheres~\cite{Dyson}. Lieb and Yngvason proved the matching
leading lower bound for repulsive interactions~\cite{LiebYngvason}.
Yau and Yin obtained the Lee--Huang--Yang upper bound for smooth
repulsive potentials~\cite{YauYin}; Basti, Cenatiempo, and Schlein
extended the upper bound to nonnegative radial compactly supported
$L^3$ interactions, with a power-saving remainder~\cite{BCSUpper}.
Fournais and Solovej proved the corresponding lower bound, first for
integrable interactions and then in a class including hard
cores~\cite{FS2020,FS2023}. The sharp hard-sphere upper bound was
subsequently obtained by Basti, Brooks, Cenatiempo, Olgiati, and
Schlein~\cite{BBCOS}. For the present class of potentials, we use the
upper bound of~\cite[Theorem~1.1]{BCSUpper} and retain the additional
positive operators in a local lower bound to recover occupation
observables.

Condensation and depletion are also understood in regimes where the
spatial and dilute limits are linked. Lieb and Seiringer proved complete
condensation for trapped gases in the Gross--Pitaevskii
limit~\cite{LS2002}. Boccato, Brennecke, Cenatiempo, and Schlein
determined the depletion in a periodic unit box with scattering length
of order $N^{-1}$~\cite[Appendix~A]{BBCS}. Their answer is a discrete
Bogoliubov momentum sum. Results of Fournais and later developments
extend condensation estimates to boxes larger than the healing
length~\cite{Fournais2021,Junge2026}. Such estimates do not directly
control the ordinary thermodynamic limit, in which the one-particle
gap closes at every fixed density. For fixed bounded interactions,
the companion~\cite{C} establishes condensation at sufficiently small
fixed densities and sufficiently low positive temperatures, and supplies
the local density tails used here.
The companion~\cite{D} establishes the depletion asymptotic for hard
spheres. The present proof adapts its separation of local energy
estimates from long-distance comparisons to bounded interactions.

Several established methods enter that adaptation. The energy analysis
uses scattering-renormalized positive squares, Neumann localization,
and Bogoliubov diagonalization, as developed in~\cite{FS2020,FS2023}
and in the low-temperature free-energy work
\cite{HHNST,FJGMOTr}. The coherent-state substitution has the standard
upper-symbol corrections of Lieb, Seiringer, and
Yngvason~\cite{LSY2005}. The probabilistic comparison averages a change
of law over routes before taking its moments. This has a direct
antecedent in the deletion-tolerance argument of Peres and
Sly~\cite[Proposition~2.1]{PeresSly}. Its geometric ingredient is the
unpredictable-path method of Benjamini, Pemantle, and
Peres~\cite{BPP}, which produces paths with exponentially decaying
intersection tails. The precise finite-volume version, with rare
encounters away from an endpoint, is supplied by
\cite[Proposition~7.1]{D}.

\subsection{From local occupations to the constant mode}

Here is the obstruction and the division of the proof. Let $\Psi$ be
the positive normalized ground function. For a partition of the torus
into cubes of side $\ell$, let $P_\ell$ be the one-particle projection
onto functions constant on each cube, and let $P_0$ project onto the
single global constant. Since $P_0\le P_\ell$, the depletion splits
exactly as
\begin{equation}\label{eq:intro-split}
 1-B_{|\Psi\rangle\langle\Psi|}
 =\langle\Psi,(1-P_\ell)_1\Psi\rangle
                    +\|(P_\ell-P_0)_1\Psi\|^2.
\end{equation}
The subscript indicates action on the first particle. The first term
measures variations inside the cubes; the second measures variations
between their mean values.

For a suitable mesoscopic choice of $\ell$, Section~\ref{sec:energy}
determines the first term at the depletion scale, together with second
and fourth moments of particle counts.
The argument follows the occupation estimates in
\cite[Section~3]{D}, but deals directly with the bounded potential.
In particular, the Neumann mirror interaction is compared with the
original interaction in expectation using a two-particle bound. This
avoids a pointwise radial-monotonicity comparison. The local lower
bound retains both a kinetic gap and positive quasiparticle terms;
these are needed to transfer the energy calculation to the one-body
observable in~\eqref{eq:intro-split}.

The second term requires information that is uniform in the total
volume. To obtain it, condition a distinguished particle to lie in a specified
cube and integrate out its position. The resulting law of the other
$N-1$ particles has density proportional to the cube average of $\Psi^2$.
Comparing these deletion laws controls the spatial variation of the
root mean squares of $\Psi$. A local comparison of deletion laws,
together with a fine-scale energy estimate, then connects those root
mean squares to the arithmetic means appearing in $P_\ell\Psi$.

Sections~\ref{sec:paths}--\ref{sec:long-palm} perform these comparisons
using ground-state Brownian paths. The construction pairs the
midpoints of paths in two copies so that deleting different endpoint
particles leaves exactly the same observed particles. For a bounded
interaction the conditional path law has Boltzmann weights. We
represent each weight by an auxiliary uniform-variable test and keep
these tests in every conditional normalization. This gives the exact
interacting law while allowing the collision estimates used in the
hard-sphere construction. The density tails of
\cite[Lemmas~4.2 and~4.3]{C} replace the deterministic count
bounds available for hard spheres.

Two different accuracies are needed. A fourth-moment comparison between
neighboring small cubes controls the passage from fine to coarse
averages. A second-moment comparison between distant cubes controls
the remaining variation across the entire torus. For the latter,
routes may have arbitrarily many steps. The estimate therefore charges
only encounters between two independently sampled routes, rather than
one fixed loss at each step. In the copy used to measure the change of law, paths near the common
endpoint remain unchanged. The route theorem makes encounters between
the changed portions rare. These conditional-law and normalization
estimates extend the local and distant comparisons of
\cite[Sections~5,~6, and~8]{D} to bounded repulsion; the hard-sphere
depletion theorem itself is not an input.

Section~\ref{sec:setup} fixes the scales and states the density-tail
inputs. Section~\ref{sec:energy} proves the local occupation and count
estimates, Section~\ref{sec:paths} constructs the conditional path law,
and Sections~\ref{sec:local-palm} and~\ref{sec:long-palm} prove the two
deletion comparisons. Section~\ref{sec:assembly} combines them through
\eqref{eq:intro-split} and restores the original units and order of
limits.
